% Part: sets-functions-relations
% Chapter: relations
% Section: reflections
%
\documentclass[../../../include/open-logic-section]{subfiles}

\begin{document}

\olfileid{sfr}{rel}{ref}
\olsection{فلسفیانہ سوچ بچار}

\olref[set]{sec} وچ اسی تعلقاں دی تعریف کجھ خاص سیٹاں دے طور تے کیتی سی۔
سانوں ذرا رک کے اک فلسفیانہ سوال پچھنا چاہیدا اے: ایہو جیہی تعریف
\emph{کر کی رہی اے}؟ ایس گل اُتے بہت شک اے کہ اسی ایہ کہنا چاہواں گے کہ
اسی مابعد الطبیعیاتی یکسانیت دے کجھ حقائق \emph{لبھ لَئے} نیں؛ مثال دے
طور تے، کہ $\Nat$ اُتے ترتیبی تعلق اوہ سیٹ $R=  \Setabs{\tuple{n,m}}{n, m \in \Nat\text{ تے }
n < m}$ \emph{نکلیا} جیہدی تعریف اسی \olref[set]{sec} وچ کیتی سی۔
ایہ گل کیوں مشکوک اے، ایہدیاں تِن وجہاں نیں۔

پہلی: \olref[set][pai]{wienerkuratowski} وچ اسی $\tuple{a, b} =
\{\{a\}, \{a, b\}\}$ دی تعریف کیتی سی۔ ایہدی تھاں ایہ تعریف ویکھو:
$\lVert a, b\rVert = \{\{b\}, \{a, b\}\} = \tuple{b,a}$۔ جدوں $a \neq b$ ہووے،
تاں $\tuple{a, b} \neq \lVert a,b\rVert$ ہوندا اے۔ پر اسی $\lVert a,b\rVert$
نوں وی ترتیب وار جوڑے دی اپنی تعریف منّ سکدے ساں، بجائے $\tuple{a,b}$ دے۔
دونویں تعریفاں اکو جِنّیاں چنگیاں طرح کم کردیاں۔ سو ہن ساڈے کول قدرتی عدداں
اُتے ترتیبی تعلق ``ہون'' دے دو اکو جِنّے چنگے امیدوار نیں، یعنی:
\begin{align*}
		R &= \Setabs{\tuple{n,m}}{n, m \in \Nat \text{ تے }n < m}\\
		S &= \Setabs{\lVert n,m\rVert}{n, m \in \Nat \text{ تے }n < m}.
\end{align*}
عنصراں نال تعیّن دی بنا اُتے $R \neq S$ اے؛ سو صاف اے کہ اوہ
\emph{دونویں} $\Nat$ اُتے ترتیبی تعلق دے نال اکو نہیں ہو سکدے۔ پر ایہ
دعویٰ کرنا کہ دراصل $R$، نہ کہ $S$ (یا ایہدے اُلٹ)، ترتیبی تعلق
\emph{اے}، محض من مرضی دی گل ہووے گی، تے ایس لئی کجھ شرمندگی والی وی۔
(ایہ سیٹ نظریاتی تحویلیت دے خلاف اوس دلیل دی اک بہت سادہ مثال اے جیہنوں
بیناسیراف نے \citeyear{Benacerraf1965} وچ مشہور کیتا سی۔ اسی ایہنوں کئی
واری مُڑ ویکھاں گے۔)

دوجی: جے اسی سمجھدے آں کہ \emph{ہر} تعلق نوں اک سیٹ دے نال اکو منّنا
چاہیدا اے، تاں خود سیٹ دا رکن ہون دا تعلق، $\in$، اک خاص سیٹ ہونا چاہیدا
اے۔ دراصل، اوہنوں سیٹ $\Setabs{\tuple{x,y}}{x \in y}$ ہونا پئے گا۔ پر کی
ایہ سیٹ موجود اے؟ رسل دے تضاد نوں سامنے رکھدیاں، ایہ کوئی معمولی دعویٰ
نہیں کہ ایہو جیہا سیٹ موجود اے۔ درحقیقت،
\oliflabeldef{cumul:::part}{سیٹاں دا اوہ نظریہ جیہڑا اسی
\olref[cumul][][]{part} وچ تیار کراں گے، ایس سیٹ دی ہوند دا \emph{انکار}
کرے گا۔\footnote{ذرا اگے دی گل کر کے وجہ ویکھ لئیے۔ خلف توں ثبوت لئی فرض
کرو کہ $I = \Setabs{\tuple{x,y}}{x \in y}$ موجود اے۔ فیر $\bigcup \bigcup I$
ہمہ گیر سیٹ اے، جیہڑا \olref[sfr][z][sep]{thm:NoUniversalSet} دے خلاف اے۔}}{سیٹ
نظریے نوں اک مسلّمی نظریے دے طور تے پوری سختی نال تیار کرنا ممکن اے، تے
اوہ نظریہ واقعی ایس سیٹ دی ہوند دا انکار کرے گا۔}
سو بھاویں کجھ تعلقاں نوں سیٹ منّ کے کم چلایا جا سکدا اے، سیٹ دا رکن ہون
دے تعلق نوں اک خاص صورت منّنا پئے گا۔

تیجی: جدوں اسی تعلقاں نوں سیٹاں دے نال ``اکو'' منّیا، تاں اسی آکھیا سی کہ
اسی $Rxy$ نوں $\tuple{x,y} \in R$ دی تھاں لکھن دی اجازت دے دیاں گے۔ ایہ
ٹھیک اے، بشرطیکہ رکن ہون دے تعلق، ``$\in$''، نوں محمول \emph{دے طور تے}
لیا جاوے۔ پر جے اسی سمجھدے آں کہ ``$\in$'' کسے خاص قسم دے سیٹ دا ناں اے،
تاں عبارت ``$\tuple{x,y} \in R$'' صرف تِن واحد اشارتی عبارتاں توں بنی اے
جیہڑیاں سیٹاں دے ناں نیں: ``$\tuple{x,y}$''، ``$\in$'' تے ``$R$''۔ تے
ناواں دی ایہو جیہی فہرست اک قضیہ بیان کرن جوگی اوہدے نالوں ودھ نہیں جِنّی
ایہ بے معنی لڑی اے: ``پیالی قلم دان میز''۔ فیر اوہی گل اے: بھاویں کجھ
تعلقاں نوں سیٹ منّ کے کم چلایا جا سکدا اے، سیٹ دا رکن ہون دے تعلق نوں اک
خاص صورت منّنا پئے گا۔ (ایہ فریگے دے تصور \emph{گھوڑا} والے تضاد دی اک
سادہ صورت تے وٹگنسٹائن دے رسل اُتے کیتے اک مشہور اعتراض نوں اکٹھا کردا اے۔)

سو ہن اسی کیہڑے نتیجے اُتے پہنچے آں؟ عدداں اُتے تعلقاں نوں سیٹ آکھن وچ
کوئی \emph{غلطی} نہیں۔ بس سانوں سمجھنا چاہیدا اے کہ ایہ گل کس بھاوَ نال
آکھی جاندی اے۔ اسی مابعد الطبیعیاتی یکسانیت دا کوئی حقیقتی دعویٰ نہیں کر
رہے۔ اسی صرف ایہ نوٹ کر رہے آں کہ کجھ سیاقاں وچ اسی (کجھ) تعلقاں نوں کجھ
خاص سیٹاں دے طور تے \emph{لے} سکدے آں (تے لے کے کم کراں گے)۔

\end{document}
